\chapter{Preliminaries}
\label{chp:prelims}


\section{Interactive Theorem Proving}


\section{Isabelle/HOL System}

\section{The Isolation Lemma}
The Isolation Lemma, originally introduced by Mulmuley, Vazirani, and Vazirani~\cite{mulmuley1987matching}, is a probabilistic principle used to ensure that, with high probability, a unique solution exists among a family of sets.
It has found widespread applications in randomized algorithms, particularly in problems involving combinatorial search, matching, and derandomization.

\vspace{0.5cm}
\noindent
\textbf{Informal Statement.}
Let $U$ be a finite set (called the \emph{universe}), and let $\mathcal{F} \subseteq \mathcal{P}(U)$ be a non-empty family of subsets of $U$.
Each element $x \in U$ is assigned a weight $w(x)$ independently and uniformly at random from $\{1, 2, \dots, N\}$ for some integer $N \ge 1$.
The total weight of a set $S \in \mathcal{F}$ is defined as"
\[
    w(S) = \sum_{x \in S} w(x)    
\]
The Isolation Lemma states that, with high probability, there is a unique set in $\mathcal{F}$ that minimizes this weight function.

\vspace{0.5cm}
\noindent
\textbf{Theorem (Isolation Lemma)~\cite{mulmuley1987matching}:}
\emph{
    Let $U$ be a finite set with $|U| = n$, and let $\mathcal{F} \subseteq \mathcal{P}(U)$ be a non-empty family of subsets.
    Suppose that each $x \in U$ is assigned an integer weight $w(x)$ drawn independently and uniformly at random from $\{1, 2, \dots, N\}$.
    Then the probability that there exists a unique set $A \in \mathcal{F}$ that minimizers $w(S)$ is at least $1 - \frac{n}{N}$.
}
\[
    \Pr[\exists! A \in \mathcal{F} \text{ minimizing } w(A)] \ge 1 - \frac{n}{N} 
\]

\vspace{0.5cm}
\noindent
\textbf{Proof Idea.}
The proof relies on analyzing the probability that \emph{two or more} sets achieving the same minimum weight.
The key technique is to consider the structure of conflicting pairs and show that their probability of collision is small.
By summing over all such pairs and bounding their contribution, one obtains the upperbound $\frac{n}{N}$ on the failure probability.

\vspace{0.5cm}
\noindent
\textbf{Proof (Joel Spencer~\cite{spencer1995probabilistic}).}
Let $U = \{e_1, e_2, \dots, e_n\}$ be a finite set of $n$ elements, and let $\mathcal{F} \subseteq \mathcal{P}(U)$ be a non-empty family of subset of $U$.
Suppose each element $x \in U$ is assigned a random weight $w(x)$ uniformly from $\{1, 2, \dots, N\}$.

For any $x \in U$, define:
\[
\alpha(x) = \min_{\substack{S \in \mathcal{F} \\ x \notin S}} w(S)- \min_{\substack{S \in \mathcal{F} \\ x \in S}} w(S \setminus \{x\})
\]
This value $\alpha(x)$ depends only on the weights of elements other than $x$ itself, since $x$ is excluded from all weight sums in the definition.
Because $w(x)$ is independently and uniformly chosen from $\{1, 2, \dots, N\}$, the probability that $w(x) = \alpha(x)$ is at most $1/N$.

Now, suppose there are two distinct sets $A, B \in \mathcal{F}$ such that both minimize the total weight $w(S) = \sum_{x \in S} w(x)$.
Let $x \in A \triangle B$ (i.e., an element in the symmetric difference). Without loss of generality, suppose $x \in A \setminus B$.
Then we have:
\[
\begin{aligned}
\alpha(x) &= \min_{\substack{S \in \mathcal{F} \\ x \notin S}} w(S)- \min_{\substack{S \in \mathcal{F} \\ x \in S}} w(S \setminus \{x\})
          &= w(B) - \left( w(A) - w(x) \right)
          &= w(x)
\end{aligned}
\]
So for this $x$, $w(x) = \alpha(x)$. The chance that this happens for any particular $x$ is at most $1/N$, and since there are at most $n$ such $x$, the total probability that some $x$ satisfies this is at most $n/N$.

Thus, with probability at least $1 - \frac{n}{N}$, no such $x$ exists, and the minimum-weight set in $\mathcal{F}$ is unique.

\vspace{0.5cm}
\noindent
\textbf{Remarks.}
\begin{itemize}
    \item The bound $1 - \frac{n}{N}$ improves as $N$ increases. In particular, for $N > 2n$, the probability of uniqueness exceeds $0.5$.
    \item The lemma is non-constructive: it guarantees uniqueness probabilistically, but not which set is unique.
    \item It is widely used in randomized parallel algorithms and in complexity theory, e.g., in the isolation of perfect matchings.
\end{itemize}

\vspace{0.5cm}
\noindent
\textbf{Commentary.}
While the above proof~\cite{spencer1995probabilistic} is concise and elegant, it omits many low-level details that are often assumed in traditional mathematical exposition.
For example, the definition of $\alpha(x)$ implicitly assumes that the minima involved are well-defined, and that comparisons among set weights are meaningful under all random assignments.
Additionally, the reasoning about probabilities (e.g., the independence and uniformity of weights) relies on intuition rather than formal measure-theoretic grounding.
These gaps, though acceptable in informal proofs, present significant challenges when one attempts to formalize such arguments in a theorem prover like Isabelle/HOL.
Thus, part of the motivation of this project is bridge this gap by giving a fully rigorous, machine-checked account of the Isolation Lemma, including all technical conditions and edge cases often left implicit in classical proofs.
In the next chapter, we will present a formalization of this result in Isabelle/HOL, making all implicit assumptions explicit and verifying each step within a constructive logical framework.

\section{Parallel Algorithm for Perfect Matching}
